% Authoring supplement: exact host bracket abstraction and a distinct later source route.
\section{Eliminating binders by combinators}
\label{sec:p07-combinatory-abstraction}

Bracket abstraction compiles a bound variable into an expression built from
application and two fixed combinators.  Its basic guarantee concerns what the
compiled expression does when it receives an argument.  The equality used to
compare unapplied expressions must also be specified.  This chapter supplies
the algorithmic content of Theorem~\ref{thm:p07-bracket-abstraction}, using the
existing S--K contractions on Sheet~\ref{sheet:sk-combinatory}.
It also gives a separate reconstruction of Selinger's atom-by-atom abstractor
and combinatory-completeness argument
\cite[Sections 5.1--5.4; Lemma 5.10; Section 5.6]{Selinger2013LambdaNotesV2}.

\subsection{Terms, replacement, and the two abstraction algorithms}
\label{sec:p07-combinatory-algorithms}

\begin{hbdefinition}{def:p07-combinatory-syntax}{The two syntactic domains}
Fix a countably infinite set of variables.  Write \(\mathcal C\) for the finite
S--K terms
\[
 P,Q ::= x\mid\mathsf S\mid\mathsf K\mid(PQ),
\]
and \(\Lambda\) for the finite lambda terms
\[
 M,N ::= x\mid(MN)\mid\lambda x.M.
\]
Application associates to the left in both domains.  Lambda terms are identified
up to renaming of bound variables; combinatory terms have no binders.  The set
\(\FV(P)\) is the set of variable leaves of \(P\), with
\(\FV(\mathsf S)=\FV(\mathsf K)=\varnothing\) and
\(\FV(PQ)=\FV(P)\cup\FV(Q)\).  For lambda terms,
\(\FV(\lambda x.M)=\FV(M)\setminus\{x\}\).

The expression \(P[Q/x]\) replaces every variable leaf \(x\) by \(Q\), without
changing \(\mathsf S\) or \(\mathsf K\).  It is ordinary tree substitution.
Lambda substitution \(M[N/x]\) is capture avoiding: at a binder different from
\(x\), first rename its variable away from \(\FV(N)\cup\{x\}\); substitution
does not pass a binder for \(x\).  Fresh variables exist because all terms are
finite.  Let \(\mathsf I:=\mathsf S\mathsf K\mathsf K\); \(\mathsf I\) is
an abbreviation, not an additional primitive.
\end{hbdefinition}

\begin{hbdefinition}{def:p07-combinatory-abstraction-algorithms}{Host and atom-by-atom abstraction}
For \(P\in\mathcal C\), the host algorithm \(\mathcal B_x(P)\) is
\[
\mathcal B_x(P):=
\begin{cases}
 \mathsf I,&P=x,\\
 \mathsf K P,&x\notin\FV(P),\\
 \mathsf S\mathcal B_x(U)\mathcal B_x(V),&P=UV,
                    \ x\in\FV(P).
\end{cases}
\]
These cases are disjoint and exhaustive.  In the last case each recursive call
has a proper subtree as input.  In particular the second case applies to an
entire application independent of \(x\), not only to an atom.

Selinger's construction is denoted here by \(\mathcal A_x\).  Its four cases are
\[
\begin{aligned}
 \mathcal A_x(x)&:=\mathsf I,\\
 \mathcal A_x(y)&:=\mathsf K y\quad(y\ne x),\\
 \mathcal A_x(c)&:=\mathsf K c\quad(c\in\{\mathsf S,\mathsf K\}),\\
 \mathcal A_x(UV)&:=\mathsf S\mathcal A_x(U)\mathcal A_x(V).
\end{aligned}
\]
In the source's polynomial language the third case applies to every coefficient
of its applicative structure.  Thus \(\mathcal A_x(yz)=
\mathsf S(\mathsf K y)(\mathsf K z)\), whereas
\(\mathcal B_x(yz)=\mathsf K(yz)\) when \(x,y,z\) are distinct.
The source writes its abstractor as \(\lambda^*x.t\); the two letters here
identify the two algorithms without treating their outputs as identical.
\end{hbdefinition}

\begin{hbdefinition}{def:p07-combinatory-relations}{Reduction and three equalities}
Write \(P\longrightarrow_{\mathsf{SK}}Q\) for one contraction of the existing
S--K sheet in any application context, and
\(P\longrightarrow_{\mathsf{SK}}^*Q\) for its reflexive transitive closure.
Its symmetric transitive reflexive closure is \(=_{\mathsf{SK}}\).
For lambda terms, \(=_\beta\) is conversion generated by capture-avoiding
beta contraction in all lambda contexts.  The relation \(=_{\beta\eta}\)
also permits \(\lambda x.Mx\longrightarrow_\eta M\) when
\(x\notin\FV(M)\).

Extensional combinatory equality \(=_{\mathrm e}\) is the least equivalence
relation on \(\mathcal C\) closed under application, containing the two S--K
equations, and closed under the fresh-variable extensionality rule on
Sheet~\ref{sheet:p07-combinatory-abstraction}.  It is a separately defined
equational theory; it adds no contraction to
\(\longrightarrow_{\mathsf{SK}}\).
\end{hbdefinition}

\begin{hbtheorem}{thm:p07-combinatory-substitution}{proved=annex:p07.combinatory.substitution}{Abstraction, free variables, and substitution}
For \(\mathcal D\in\{\mathcal A,\mathcal B\}\), every term \(P\), and every
variable \(x\), the recursive algorithm terminates and
\[
 \FV(\mathcal D_x(P))=\FV(P)\setminus\{x\},\qquad
 \mathcal D_x(P)Q\longrightarrow_{\mathsf{SK}}^*P[Q/x]
 \quad(Q\in\mathcal C).
\]
If \(y\ne x\) and \(x\notin\FV(Q)\), the host algorithm satisfies literal
equality of combinatory trees:
\[
 \mathcal B_x(P)[Q/y]=\mathcal B_x(P[Q/y]).
\]
For \(\mathcal A\) the corresponding general statement is
\(\mathcal A_x(P)[Q/y]=_{\mathrm e}\mathcal A_x(P[Q/y])\); it is literal
when \(Q\) is an atom.  For example, if \(x,y,u,v\) are distinct, substitution
of \(uv\) into \(\mathcal A_x(y)=\mathsf Ky\) gives \(\mathsf K(uv)\),
whereas recomputation gives \(\mathsf S(\mathsf Ku)(\mathsf Kv)\).
If \(z\ne x\) and \(z\notin\FV(P)\), then
\(\mathcal D_x(P)=\mathcal D_z(P[z/x])\).
For distinct \(x_1,\ldots,x_n\), simultaneous substitution is implemented by
\[
 (\mathcal D_{x_1}(\mathcal D_{x_2}(\cdots
 \mathcal D_{x_n}(P)\cdots)))Q_1\cdots Q_n
 \longrightarrow_{\mathsf{SK}}^*P[Q_1/x_1,\ldots,Q_n/x_n]
\]
provided none of the \(x_i\) occurs in any \(Q_j\).
This last freshness condition distinguishes simultaneous replacement from
successive replacement that might also act inside previously inserted arguments.
\end{hbtheorem}

\subsection{Translation to lambda terms and simulation by application}
\label{sec:p07-combinatory-translations}

\begin{hbdefinition}{def:p07-combinatory-translations}{Expansion and compilation}
The expansion \(E:\mathcal C\to\Lambda\) is
\[
\begin{aligned}
 E(x)&=x,& E(PQ)&=E(P)E(Q),\\
 E(\mathsf K)&=\lambda a.\lambda b.a,&
 E(\mathsf S)&=\lambda a.\lambda b.\lambda c.ac(bc).
\end{aligned}
\]
Bound names in these closed expressions may always be chosen fresh.
For \(\mathcal D\in\{\mathcal A,\mathcal B\}\), compilation
\(T_{\mathcal D}:\Lambda\to\mathcal C\) is
\[
 T_{\mathcal D}(x)=x,\quad
 T_{\mathcal D}(MN)=T_{\mathcal D}(M)T_{\mathcal D}(N),\quad
 T_{\mathcal D}(\lambda x.M)=\mathcal D_x(T_{\mathcal D}(M)).
\]
It removes inner binders before applying the outer abstraction algorithm.
Both maps act on finite syntax.  Their effective computation requires neither
normalization of the input nor decidability of conversion.
\end{hbdefinition}

\begin{hbtheorem}{thm:p07-combinatory-translation}{proved=annex:p07.combinatory.translation}{Capture facts and beta-correct expansion}
For either algorithm, compilation is independent of bound-variable renaming
and preserves the exact free-variable set.  The host compiler commutes
literally with capture-avoiding substitution:
\[
 T_{\mathcal B}(M[N/x])=
 T_{\mathcal B}(M)[T_{\mathcal B}(N)/x].
\]
For \(T_{\mathcal A}\), this is an \(=_{\mathrm e}\) equality in general
and a literal equality when \(N\) is a variable.
Expansion also preserves free variables and commutes with substitution up to
renaming bound variables.  Moreover,
\[
 E(\mathcal D_x(P))=_\beta\lambda x.E(P),\qquad
 E(T_{\mathcal D}(M))=_\beta M.
\]
Every weak S--K step expands to a nonempty finite beta-reduction sequence.
A beta contraction at the root has the finite S--K simulation
\[
 T_{\mathcal D}((\lambda x.M)N)
 \longrightarrow_{\mathsf{SK}}^*
 T_{\mathcal D}(M)[T_{\mathcal D}(N)/x].
\]
For the host compiler this target is exactly \(T_{\mathcal B}(M[N/x])\).
For the atom-by-atom compiler it is extensionally equal to the recompilation;
literal recompilation is not asserted.
The same assertion holds in contexts generated by
\(C[-]::=[-]\mid C[-]M\mid M C[-]\).
These contexts do not place the contracted redex beneath a lambda binder.
\end{hbtheorem}

\begin{hbdefinition}{def:p07-combinatory-beta-presentation}{Representing lambda conversion and numerical outputs}
Define \(P=_{\lambda\beta}Q\) to mean \(E(P)=_\beta E(Q)\).
This equality is a transported beta theory, distinct from weak S--K conversion.
With this explicit equality, compilation preserves and reflects beta conversion:
\[
 M=_\beta N\quad\Longleftrightarrow\quad
 T_{\mathcal D}(M)=_{\lambda\beta}T_{\mathcal D}(N).
\]
This equivalence follows in both directions from
Theorem~\ref{thm:p07-combinatory-translation} and transitivity of beta conversion.
It permits a precise representation statement for the Church numerals of
Definition~\ref{def:p07-church-numeral}: write \(\nu_n\) for that lambda term
and \(c_n=T_{\mathcal D}(\nu_n)\).  A closed lambda term \(F\) has a Church
numeral \(\nu_m\) as the beta-normal output on inputs \(n_1,\ldots,n_k\)
exactly when
\[
 T_{\mathcal D}(F)c_{n_1}\cdots c_{n_k}
       =_{\lambda\beta}c_m.
\]
Indeed expansion of the left side is beta-convertible to
\(F\nu_{n_1}\cdots\nu_{n_k}\); beta confluence
(Theorem~\ref{thm:p07-church-rosser}) converts equality with the normal term
\(\nu_m\) into a reduction to it.  Undefined numerical inputs, characterized
by absence of any numeral normal output, are preserved in this same sense.
This statement specifies the observation used; it does not assert that weak
S--K reduction reaches the particular tree \(c_m\).
\end{hbdefinition}

\subsection{Why weak conversion does not supply abstraction congruence}
\label{sec:p07-combinatory-confluence}

\begin{hbdefinition}{def:p07-combinatory-parallel}{Auxiliary parallel S--K reduction}
The relation \(\Rightarrow_{\mathsf{SK}}\) on binder-free S--K terms is
generated by the four parallel-reduction clauses on
Sheet~\ref{sheet:p07-combinatory-abstraction}. Application may reduce both
operands in parallel. A root S contraction uses the same selected reduct
of its third operand in both resulting copies. This auxiliary relation is
distinct from one weak contraction and is used in the confluence proof.
\end{hbdefinition}
\showcalculus{p07-combinatory-abstraction}

\begin{hbtheorem}{thm:p07-combinatory-confluence}{proved=annex:p07.combinatory.confluence}{Confluence of the two S--K contractions}
Weak S--K reduction is confluent: if
\(P\longrightarrow_{\mathsf{SK}}^*Q\) and
\(P\longrightarrow_{\mathsf{SK}}^*R\), there is \(U\) to which both \(Q\)
and \(R\) reduce.  Consequently two weak S--K normal forms are
\(=_{\mathsf{SK}}\)-equal only when they are the same tree.
\end{hbtheorem}

\begin{hbtheorem}{thm:p07-combinatory-weak-boundary}{proved=annex:p07.combinatory.weak-boundary}{Beta equality need not compile to weak conversion}
For the host compiler, the beta-equal terms
\[
 M=\lambda x.x,\qquad N=\lambda x.(\lambda y.y)x
\]
compile to the distinct S--K normal forms
\[
 T_{\mathcal B}(M)=\mathsf I,\qquad
 T_{\mathcal B}(N)=\mathsf S(\mathsf K\mathsf I)\mathsf I.
\]
They are not weakly convertible.  Nevertheless, after application to any \(Q\),
both reduce to \(Q\).  Thus neither beta conversion beneath binders nor equality
of all application results may be silently replaced by \(=_{\mathsf{SK}}\).
\end{hbtheorem}

\subsection{Extensional correspondence and an eta optimization}
\label{sec:p07-combinatory-extensional}

\begin{hbtheorem}{thm:p07-combinatory-extensional}{proved=annex:p07.combinatory.extensional}{Exact correspondence for the extensional theory}
For \(\mathcal D\in\{\mathcal A,\mathcal B\}\), bracket abstraction respects
\(=_{\mathrm e}\): if \(P=_{\mathrm e}Q\), then
\(\mathcal D_x(P)=_{\mathrm e}\mathcal D_x(Q)\).
For all lambda terms and combinatory terms,
\[
\begin{split}
 M=_{\beta\eta}N&\quad\Longleftrightarrow\quad
 T_{\mathcal D}(M)=_{\mathrm e}T_{\mathcal D}(N),\\
 P=_{\mathrm e}Q&\quad\Longleftrightarrow\quad
 E(P)=_{\beta\eta}E(Q),\\
 T_{\mathcal D}(E(P))&=_{\mathrm e}P.
\end{split}
\]
In particular \(\mathcal A_x(P)=_{\mathrm e}\mathcal B_x(P)\).
These are equations in the named theories, rather than identities of trees or
claims of reduction in a specified direction.
\end{hbtheorem}

\begin{hbdefinition}{def:p07-combinatory-eta-algorithm}{An explicitly different eta-optimized algorithm}
Define \(\mathcal H_x(P)\) by the following ordered cases: use \(\mathsf I\)
if \(P=x\); use \(\mathsf K P\) if \(x\notin\FV(P)\); use \(U\) if
\(P=Ux\) and \(x\notin\FV(U)\); otherwise, for \(P=UV\), use
\(\mathsf S\mathcal H_x(U)\mathcal H_x(V)\).
Define \(T_{\mathcal H}\) by the same compilation recursion with
\(\mathcal H\) in place of \(\mathcal D\).
The fresh-variable condition in the third case is indispensable.
\end{hbdefinition}

\begin{hbtheorem}{thm:p07-combinatory-eta-boundary}{proved=annex:p07.combinatory.eta-boundary}{What the eta optimization preserves}
For all \(P,Q,x\),
\[
\FV(\mathcal H_x(P))=\FV(P)\setminus\{x\},\qquad
\mathcal H_x(P)Q\longrightarrow_{\mathsf{SK}}^*P[Q/x],\qquad
E(\mathcal H_x(P))=_{\beta\eta}\lambda x.E(P).
\]
Consequently \(E(T_{\mathcal H}(M))=_{\beta\eta}M\).
The beta-only analogue fails: for distinct \(x,f\),
\(\mathcal H_x(fx)=f\), but \(f\ne_\beta\lambda x.fx\).
If the freshness condition is dropped, the proposed simplification
\(\mathcal H_x(xx)=x\) fails both variable elimination and the application
property.
\end{hbtheorem}

\subsection{Representation in applicative structures}
\label{sec:p07-combinatory-algebras}

\begin{hbdefinition}{def:p07-combinatory-algebra}{Polynomial operations and combinatory completeness}
An applicative structure is a nonempty set \(A\) with a total binary operation
\(\cdot:A\times A\to A\), written by juxtaposition.  Application need not be
associative.  A polynomial over \(A\) is a finite tree built from variables,
named coefficients \(a\in A\), and application.  An assignment \(\rho\) of
variables to elements evaluates that tree by interpreting its coefficients as
the named elements.  Write its value as \(\llbracket p\rrbracket_\rho\).

The structure is combinatorially complete when every polynomial with variables
among distinct \(x_1,\ldots,x_n\), for every \(n\geq0\), has a representing
element \(a\in A\) satisfying
\[
 a b_1\cdots b_n=
 \llbracket p\rrbracket_{[x_1\mapsto b_1,\ldots,x_n\mapsto b_n]}
 \quad\text{for all }b_1,\ldots,b_n\in A.
\]
For \(n=0\) this says that a closed polynomial has its evaluated value as
representative.  This is a statement about polynomial operations, not about
every set-theoretic function \(A^n\to A\).  Representatives need not be unique.
A combinatory algebra specifies elements \(s,k\in A\) such that
\(sabc=(ac)(bc)\) and \(kab=a\) for every \(a,b,c\in A\).
\end{hbdefinition}

\begin{hbtheorem}{thm:p07-combinatory-completeness}{proved=annex:p07.combinatory.completeness}{S and K characterize polynomial representability}
An applicative structure is combinatorially complete if and only if it has
elements \(s,k\) making it a combinatory algebra.  Either abstraction algorithm,
with coefficients interpreted as constants and \(\mathsf S,\mathsf K\)
interpreted as \(s,k\), constructs a representative of every polynomial.
It also satisfies the assignment-sensitive identity
\[
 \llbracket\mathcal D_x(p)\rrbracket_\rho\,b
       =\llbracket p\rrbracket_{\rho[x\mapsto b]}.
\]
If the structure is extensional, meaning that
\((\forall c\in A)\,ac=bc\) implies \(a=b\), the interpretation also respects
\(=_{\mathrm e}\) and hence beta-eta equality of compiled lambda terms.
\end{hbtheorem}

\begin{hbexample}{def:p07-combinatory-examples}{Arguments that select, duplicate, and preserve a free variable}
For arbitrary terms \(U,V\), \(\mathsf KUV\longrightarrow_{\mathsf{SK}}U\),
and \((\mathsf K\mathsf I)UV\longrightarrow_{\mathsf{SK}}^*V\).
Also \(\mathcal B_x(xx)=\mathsf S\mathsf I\mathsf I\), and application to
\(U\) reduces to \(UU\).
For distinct variables \(x,y\),
\[
 \mathcal B_x(yx)=\mathsf S(\mathsf K y)\mathsf I,
 \qquad
 \mathsf S(\mathsf K y)\mathsf I U
 \longrightarrow_{\mathsf{SK}}(\mathsf K yU)(\mathsf I U)
 \longrightarrow_{\mathsf{SK}}^*yU.
\]
The variable \(y\) remains free.  Capture avoidance is visible in
\((\lambda x.\lambda y.x)y\): rename the inner binder to \(z\) before
substitution, obtaining \(\lambda z.y\), whose compilation is
\(\mathsf K y\).  Treating the result as \(\lambda y.y\) would instead compile
to \(\mathsf I\) and would change the free-variable set and behavior.
\end{hbexample}

\subsection{Finite size bounds and use of the translation}
\label{sec:p07-combinatory-size}

\begin{hbdefinition}{def:p07-combinatory-size}{Tree size and binder count}
Let \(s(P)\) count variable/constant leaves and application nodes in a combinatory
tree; thus \(s(x)=s(\mathsf S)=s(\mathsf K)=1\) and
\(s(PQ)=1+s(P)+s(Q)\).  For lambda terms count each lambda binder as one more
node and use the same leaf/application counts.  Let \(b(M)\) be the number of
lambda binders.  The abbreviation \(\mathsf I\) is expanded to its five-node
S--K tree when measuring size.  These are tree counts, not a claim about shared
graphs, printed variable-name lengths, or reduction time.
\end{hbdefinition}

\begin{hbtheorem}{thm:p07-combinatory-size}{proved=annex:p07.combinatory.size}{Explicit bounds for the stated algorithms}
For \(\mathcal D\in\{\mathcal A,\mathcal B,\mathcal H\}\),
\[
 s(\mathcal D_x(P))\leq4s(P)+1\leq5s(P),\qquad
 s(T_{\mathcal D}(M))\leq5^{b(M)}s(M),\qquad
 s(E(P))\leq10s(P).
\]
The first bound is attained by every application tree whose leaves are all
\(x\), for \(\mathcal A\) and \(\mathcal B\).
A single abstraction can be implemented in time linear in its input and output
tree size after marking, in one bottom-up pass, whether each subtree contains
\(x\).  The displayed bound for full compilation is an upper bound; no
polynomial bound or matching exponential lower bound is asserted.
\end{hbtheorem}

\begin{hbmethod}{def:p07-combinatory-lookup}{Choosing the required guarantee}
For an argument-passing computation, use the substitution theorem and identify
the application context.  To reconstruct a lambda expression modulo beta,
use \(T_{\mathcal A}\) or \(T_{\mathcal B}\) followed by \(E\), or specify the
transported relation \(=_{\lambda\beta}\).  To compare extensional functions,
state \(=_{\mathrm e}\), with its fresh-variable rule and the corresponding
beta-eta theorem.  An eta optimization uses the guarantees of
Theorem~\ref{thm:p07-combinatory-eta-boundary}.  When a compiled function is
compared before being applied, the equality choice is part of the claim.
\end{hbmethod}
\index{bracket abstraction!substitution theorem}
\index{combinatory logic!extensional equality}
